\chapter{Course Conclusion and Further Topics\texorpdfstring{$^{(*)}$}{(*)}}
\label{chap:conclusion}

The course progressed from the \emph{Linear Quadratic Regulator (LQR)} through \emph{Model Predictive Control (MPC)} with explicit constraint handling, to \emph{State Estimation} and \emph{Disturbance Modelling} for offset-free tracking in the presence of noise and uncertainty.

The structure presented here---\emph{MPC + Estimator + Target Calculator}---is the standard architecture for linear MPC in process control. We omitted several advanced topics; this chapter provides pointers for further study.

\section{Rigorous Stability Analysis}
In the basic MPC formulation without a terminal set constraint, the \emph{Terminal Cost} ($x_N^\top P x_N$) encourages stability but does not mathematically \textit{guarantee} it for constrained systems. (Chapter~\ref{chap:linear_mpc}, Section~\ref{sec:mpc_stability_theory} provides the rigorous proof when the terminal set \emph{is} included.)

Rigorous MPC theory relies on Lyapunov analysis. To prove stability, one must often include a \emph{Terminal Set Constraint} ($x_N \in \mathcal{X}_f$) in the optimisation. This forces the state to land in a \emph{safe haven} (a Positive Invariant Set, meaning that once the trajectory enters the terminal set, it remains inside) at the end of the horizon, ensuring recursive feasibility.

The fundamental reference for this topic is Mayne et al.~\cite{Mayne2000}.

\section{Robust MPC}
Robust MPC, specifically the Tube MPC approach, is covered in detail in Chapter~\ref{chap:tube_mpc}. The reader is referred there for the full development including Minkowski sums, Pontryagin differences, minimal robust positively invariant sets, and constraint tightening.

\section{Nonlinear MPC (NMPC)}
Nonlinear MPC is covered in detail in Chapter~\ref{chap:nmpc}. The reader is referred there for direct transcription methods, the CasADi/IPOPT workflow, successive convexification (SCvx), GuSTO, a powered rocket landing application, and Real-Time Iteration.

\section{Real-Time Implementation and Embedded MPC}
In our examples, we used solvers such as MATLAB's \texttt{quadprog} or the open-source solver IPOPT (accessed via CasADi or other interfaces), which are sufficient for offline simulations but insufficient for millisecond-level control on microcontrollers or DSPs. Deploying MPC on embedded hardware requires specialised techniques:
\begin{itemize}
	\item \textbf{Explicit MPC (Multiparametric Programming)~\cite{Bemporad2002}:} For small systems, we can pre-calculate the solution offline. The result is a lookup table of affine gains, one per polyhedral region of the state space. On the chip, the ``optimisation'' reduces to \texttt{if-else} checks to find the correct region, requiring microseconds to execute. Section~\ref{sec:explicit_mpc} develops this approach with worked examples and MPT3 code.
	\item \textbf{Code Generation:} Tools like \href{https://acado.github.io/}{ACADO}, \href{https://cvxgen.com/}{CVXGEN}, or \href{https://osqp.org/}{OSQP} generate highly optimised, library-free C-code specifically tailored to the structure of the MPC problem, allowing convex MPC to run on standard automotive ECUs.
\end{itemize}

The fundamental reference for explicit MPC is Bemporad et al.~\cite{Bemporad2002}.

\section{Data-Driven and Learning-Based MPC}
Data-driven and learning-based MPC methods are covered in detail in Chapter~\ref{chap:learning_mpc}. The reader is referred there for GP-MPC, Koopman MPC, DeePC (including Willems' Fundamental Lemma), and Learning MPC, with full derivations, worked examples, and MATLAB/YALMIP implementations.

